\scanpage{246}
\mspage{120}
% archive sheet number: 245

\uline{Corollaire 9.2}. \deleted{\unsure{Supposons|Réciproquement}} Sous l'hyp. Loc $(1,4)$, soit $f:X\unsure{\to|-}Y$, \deleted{\unsure{soit|et}}\inserted{\unsure{$\in$|une}} flèche de $\mathrm{Cat}$, et $Z$ objet de $\mathrm{Cat}$. Alors $f\in W$ ssi $f\amalg\mathrm{id}_Z:X\amalg Z-Y\amalg Z$ est $\in W$.

Immédiat sur le carré cocartésien, où \deleted{\unsure{est imm. ouv.}} $\mathrm{inj}_0^X$ est imm. ouv. de Dwyer

\diagram{\begin{tikzcd}[ampersand replacement=\&,column sep=large]
X \arrow[r,"f",no head] \arrow[d,"\mathrm{inj}_0^X"'] \& Y \arrow[d,"\mathrm{inj}_0^Y"] \\
X\amalg Z \arrow[r,"{f'=f\amalg\mathrm{id}_Z}"',no head] \& Y\amalg Z
\end{tikzcd}}

qui montre que $f\in W\unsure{\iff|\Longrightarrow} f'\in W$.

\unsure{:} En itérant ce résultat, on trouve, si
\[
 f_1:X_1-Y_1,\quad f_2:X_2-Y_2
\]
sont donnés, alors
\[
 f_1\in W,\ f_2\in W\Longrightarrow f_1\amalg f_2\in W.
\]
en regardant le diagramme \deleted{\unsure{bi}}\inserted{de deux carrés} cocartésiens : \deleted{\unsure{d'immersions ouvertes}}

\diagram{\begin{tikzcd}[ampersand replacement=\&,column sep=huge,row sep=large]
X_1 \arrow[r,"f_1"] \arrow[d] \& Y_1 \arrow[d] \\
X_1\amalg X_2 \arrow[r,"{\underbrace{f'_1=f_1\amalg\mathrm{id}_{X_2}}}"] \arrow[d,"{f'_2=\mathrm{id}_{X_1}\amalg f_2}"] \& Y_1\amalg X_2 \arrow[d,"{g_2=\mathrm{id}_{Y_1}\amalg f_2}"] \\
X_1\amalg Y_2 \arrow[r,"{g_1=\deleted{\unsure{f}}f_1\amalg\mathrm{id}_{Y_2}}"'] \& Y_1\amalg Y_2
\end{tikzcd}}

Si $f_1$ et $f_2$ sont dans $W$, il en est de $=$ de $f'_1,f'_2$, \deleted{\unsure{par 9.2}}\inserted{$g_1,g_2$}, \deleted{et si on prouve que}\inserted{donc aussi de $f_1\amalg f_2=g_2f'_1$.}

\deleted{chaque carré est $W$-cocartésien, on en conclut \inserted{\unsure{(cf juste hyp. sur $f_1,f_2$,|sans hyp. sur $f_1,f_2$,} par Loc $(4)$)} que $g_1,g_2\in W$, donc aussi $f_1\amalg f_2=g_2f'_1\in W$. (Mais je ne vois \unsure{aucun|aucune}}
